\section{Conclusions}

Chapter~1 framed this thesis as an end-to-end engineering pipeline (from a raw optical video archive to an uncertainty-aware design-triage tool) for marine multi-fuel pilot injection screening. Its central contribution is a censoring-aware measurement-to-physics analysis: finite-FOV observations are converted into auditable trajectory targets, and those targets support a time-resolved empirical test of pressure scaling before they are used by the surrogate. The validation evidence and its interpretation are presented in Chapter~6; this chapter distils the resulting contributions and main findings.

\subsection{Novelty and Contributions}

The primary contributions of this thesis are:

\begin{enumerate}
  \item \textbf{Censoring-aware target construction:} An observation audit and reduced-order target chain that distinguishes finite-FOV and density-loss truncation from physical spray saturation. The SGQR continuation is used as an explicit shape hypothesis, with raw/teacher hybrid supervision limiting its influence where observations remain reliable.
  \item \textbf{Time-resolved empirical pressure scaling:} A reproducible, condition-clustered analysis showing that the pooled effective \(\Delta P\) exponent is not constant over the observable short-pilot-injection interval. It is approximately 0.45--0.69 in the high-support \SIrange{0.3}{0.7}{ms} window and moves towards the H--A post-breakup value as statistical support weakens; conditional ambient-pressure dependence also remains measurable.
  \item \textbf{Uncertainty-aware surrogate and screening implementation:} A multi-stage heteroscedastic surrogate and geometric interface that convert the two findings above into time-resolved wall-interception scores ($P_{\mathrm{hit}}(t)$ and cumulative exposure \(E\)) for design triage ahead of CFD.
  \item \textbf{Engineering pipeline:} An automated, CUDA-accelerated video-processing and checkpointing workflow capable of ingesting raw multi-hole Mie scattering archives (processing 2{,}992 recordings in approximately 72 minutes) into auditable plume-wise temporal observables.
\end{enumerate}

\subsection{Main Findings}

First, finite field-of-view censoring is too prevalent to treat boundary-reaching trajectories as physically saturated. The \(\mathrm{SGQR}\) formulation provides a robust, smooth target representation for this setting, while remaining a shape-based continuation hypothesis rather than a validated physical spray law.

Second, the archive-level pressure scaling is time-dependent over the observed short-pilot-injection window. The early-to-mid transient retains a Bernoulli-like imprint before the pooled exponent moves towards the classical post-breakup asymptote. The late-time transition cannot be assigned a unique time because censoring protocol, represented-condition count, and regression support change together. The result motivates the surrogate's amplitude scaling and pressure residual features without claiming a universal law or an identified internal-nozzle mechanism.

Third, the staged surrogate substantially improves penetration prediction over the refitted empirical correlations within the primary evaluation domain. The MLP and SVGP results reveal a genuine architectural trade-off rather than an unqualified winner: the SVGP gives slightly stronger point prediction, whereas the MLP provides the output structure, staged distillation pathway, and more conservative in-domain coverage required by the proposed screening interface.

Fourth, transfer is useful but bounded within the evaluated nozzle design lineage. The failure on Nozzle~0 shows that an injector-population shift combined with operating-condition extrapolation cannot be treated as routine interpolation and requires new supporting data or model adaptation.

\subsection{Scope}

These conclusions remain limited to the room-temperature, non-reacting, constant-volume-vessel conditions and the Mie-defined liquid-envelope observables represented in the archive. The downstream wall-interception score is a simplified geometric proximity proxy, not a phase-resolved prediction of harmful liquid impingement. Wide-FOV continuation, physical wall-contact ground truth, and clustered confidence intervals therefore remain the principal validation needs identified in Chapter~6.

\subsection{Closing Statement}

By combining censoring-aware measurement interpretation with time-resolved empirical scaling, this thesis shows how archival optical Mie scattering recordings can yield both a defensible physical diagnostic and an uncertainty-aware pre-screening tool. The three-stage learner is one implementation of that evidence chain rather than its central scientific result. The resulting framework can narrow candidate nozzle geometries and operating conditions prior to resource-intensive CFD simulations and engine validation trials.
